
\medskip\noindent\textit{Continuation of the proof of \SGARefLink{sga3:xvii:proposition:4-3-1}{4.3.1}.}

\medskip\noindent
\textit{(i)\(\Rightarrow\)\SGARefLink{sga3:xvii:proposition:4:3:1:item:v}{(v)}.}\quad
This follows from \SGARefLink{sga3:xvii:corollary:3:7}{3.7}.

\medskip\noindent
\textit{(v)\(\Rightarrow\)\SGARefLink{sga3:xvii:proposition:4:3:1:item:vi}{(vi)}.}\quad
Indeed, since
\[
  (\operatorname{ad}X)^{p^r}
    =\operatorname{ad}\bigl(X^{(p^r)}\bigr)
  \qquad\text{(Exposé VII\textsubscript{A}, 5.2),}
\]
\(\operatorname{ad}X\) is nilpotent when \(\mathfrak g\) is
unipotent; hence \(\mathfrak g\) is nilpotent by Engel's theorem.
Moreover, if \(\mathfrak g\) is unipotent, then its center plainly
is as well, and the reductive part of that center is therefore
trivial (\SGARefLink{sga3:xvii:lemma:4-2-2}{4.2.2}).

\medskip\noindent
\textit{(vi)\(\Rightarrow\)(iv).}\quad
Indeed, if \(G_{\overline{k}}\) contained a subgroup isomorphic to
\(\boldsymbol\mu_p\), its Lie algebra would contain a nonzero
element \(X\) such that
\[
  X^{(p)}=X
  \qquad\text{(Appendix II, 2.1),}
\]
and hence \(X^{(p^r)}=X\) for every \(r>0\).
Since \(\operatorname{ad}X\) is nilpotent because \(\mathfrak g\)
is nilpotent, and
\[
  (\operatorname{ad}X)^{p^r}
    =\operatorname{ad}\bigl(X^{(p^r)}\bigr),
\]
we necessarily have \(\operatorname{ad}X=0\).  Thus \(X\) belongs
to the reductive part of the center of
\(\mathfrak g_{\overline{k}}\), contradicting (vi).

\medskip\noindent
\textit{(i)\(\Rightarrow\)(vi ter).}\quad
This follows from \SGARefLink{sga3:VIIA:2.4:item:i}{2.4(i)} and \SGARefLink{sga3:xvii:theorem:3:5:item:i}{3.5(i)}\(\Rightarrow\)(iii).

\medskip\noindent
\textit{(vi ter)\(\Rightarrow\)(vi bis).}\quad
Indeed, if \(G\) is nilpotent, then so is its subgroup
\({}_{F}G\).  Appendix II, \SGARefLink{sga3:ii:subsection:2-2}{2.2} shows that \({}_{F}G\) is
nilpotent if and only if
\[
  \operatorname{Lie}({}_{F}G)=\operatorname{Lie}G
  \qquad\text{(Exposé VII\textsubscript{A})}
\]
is nilpotent.

\medskip\noindent
\textit{(vi bis)\(\Rightarrow\)(vi).}\quad
Let \(Z\) be the identity component of the center of \(G\), and
let \(\mathfrak r\) be the reductive part of the center of
\(\mathfrak g\).  We must prove that \(\mathfrak r=0\).
It is immediate that \(\mathfrak r\) is a characteristic
\(p\)-Lie subalgebra of
\[
  \mathfrak g=\operatorname{Lie}({}_{F}G);
\]
that is, it is stable under the functor
\(\operatorname{Aut}_{p\text{-Lie}}(\mathfrak g)\).
Consequently \(\mathfrak r\) is the Lie algebra of a
characteristic radicial subgroup \(R\) of \({}_{F}G\)
(Appendix II, \SGARefLink{sga3:ii:subsection:2-2}{2.2}).  The last assertion of \SGARefLink{sga3:xvii:lemma:4-2-2}{4.2.2} and Appendix II,
\SGARefLink{sga3:xvii:proposition:2:1}{2.1} show moreover that \(R\) is a form of
\((\boldsymbol\mu_p)^r\).  Since \(R\) is characteristic in
\({}_{F}G\), which is itself characteristic in \(G\)
(Appendix II, 1), it is invariant in \(G\), and hence central
because \(G\) is connected (\SGARefLink{sga3:ix:corollary:5:5}{Exposé IX, 5.5}).  By (vi bis),
\(R=0\), and therefore \(\mathfrak r=0\).

\SGASubsection{4.4}{Étale Groups}
\SGARefTarget{sga3:xvii:subsection:4:4}

The following proposition is an easy consequence of the Sylow
theorems and of the structure of finite \(q\)-groups
(cf.\ J.-P. Serre, \emph{Corps locaux}, Chapter IX, \S1).

\SGAStatement{Proposition 4.4.1}{Let \(G\) be a finite étale
algebraic group over an algebraically closed field.  Then \(G\)
is unipotent if and only if, for every prime number \(q\) distinct
from the characteristic \(p\) of the field, \(G\) contains no
subgroup isomorphic to \(\boldsymbol\mu_q\).}
\SGARefTarget{sga3:xvii:proposition:4:4:1}

\SGASubsection{4.5}{Abelian Varieties}
\SGARefTarget{sga3:xvii:subsection:4:5}

Let \(G\) be an abelian variety over an algebraically closed field
\(k\).  The following conditions are equivalent:
\begin{enumerate}
\item[(i)] \SGARefTarget{sga3:xvii:subsection:4:5:item:i}
  \(G\) is unipotent.
\item[(ii)] \SGARefTarget{sga3:xvii:subsection:4:5:item:ii}
  \(G=0\).
\item[(iii)] \SGARefTarget{sga3:xvii:subsection:4:5:item:iii}
  There exists an integer \(n\), relatively prime to the
  characteristic \(p\) of \(k\), such that \(G\) contains no
  subgroup isomorphic to \(\boldsymbol\mu_n\).
\end{enumerate}
